\section*{Capítulo \ref{ch:b2:palaeoenvironments} --- Fósiles y paleoambientes}

\begin{solution}{exo:b2:palaeoenvironments:1}
(a) $112/135 = 0.830$ y $112/150 = 0.747$. (b) Las conchas muertas del fondo se acumulan durante años
o más: incluyen especies que vivieron allí antes, o raramente, o que fueron transportadas, de modo que la
\omterm{def:b2:palaeoenvironments:assemblages}{asociación de muerte} promedia muchas generaciones de la comunidad. (c) Los grupos de cuerpo blando, en
particular los detritívoros y los carnívoros como los gusanos: sin esqueleto, se descomponen antes del enterramiento.
\end{solution}

\begin{solution}{exo:b2:palaeoenvironments:2}
(a) $5 \times 165 = \qty{825}{kyr}$, unos \qty{0.8}{Myr}, dentro de los 0.7 a 1 millón de años.
(b) $183.73 - 182.77 = \qty{0.96}{Myr}$, y $0.96/0.165 = 5.8$: unos seis horizontes. (c) Un
linaje cambia en mucho menos de un millón de años; solo los horizontes separan las formas
sucesivas, mientras que una zona las agrupa.
\end{solution}

\begin{solution}{exo:b2:palaeoenvironments:3}
(a) \omterm{def:b2:palaeoenvironments:ichnofacies}{Icnofacies} de Nereites: el fondo de mar profundo de un sistema turbidítico, más allá del borde de la plataforma, donde
el alimento escasea. (b) \omterm{def:b2:palaeoenvironments:ichnofacies}{Icnofacies} de Skolithos: zona intermareal o playa sumergida. (c) \omterm{def:b2:palaeoenvironments:ichnofacies}{Icnofacies} de Cruziana:
plataforma externa, entre los niveles de base del oleaje de buen tiempo y de tormenta. (d) \omterm{def:b2:palaeoenvironments:ichnofacies}{Icnofacies} de Zoophycos: en el
Mesozoico, la plataforma inferior o el talud, por debajo del nivel de base del oleaje de tormenta.
\end{solution}

\begin{solution}{exo:b2:palaeoenvironments:4}
(a) La capa A contiene grupos estenohalinos: agua de mar normal. La capa B, pobre en especies pero rica en
individuos, con restos vegetales y ninguno de los grupos estenohalinos, sugiere agua salobre, por
ejemplo un estuario. (b) Los grupos ausentes son varios, comunes y bien conservados en capas marinas
de la misma edad: su ausencia conjunta es una pauta, mientras que un fósil puede faltar por azar o
por mala preservación. (c) El sedimento (canales mareales o fluviales, láminas de fango, raíces), los
\omterm{def:b2:palaeoenvironments:fossil-kinds}{icnofósiles} (una \omterm{def:b2:palaeoenvironments:ichnofacies}{icnofacies} de Cruziana con pocas especies), y el paso lateral de la capa B
a depósitos marinos o fluviales.
\end{solution}

\begin{solution}{exo:b2:palaeoenvironments:5}
(a) $P = 1/2$ cuando $kh/s = \ln 2$: $s = 0.05 \times 0.1/0.693 = \qty{0.0072}{m.yr^{-1}}$,
\qty{7.2}{mm} al año. (b) Reducir $s$ a la mitad duplica el exponente: $P = (1/2)^2 = 0.25$, la mitad.
(c) $k$ se reduce a la mitad: $P = e^{-2.5} = 0.082$ en lugar de $e^{-5} = 0.0067$, $e^{2.5} = 12$ veces más
probable. (d) Las conchas gruesas y robustas están sobrerrepresentadas, y tanto más cuanto más lento es el enterramiento.
\end{solution}

\begin{solution}{exo:b2:palaeoenvironments:6}
(a) La cota superior es $183.73 + 0.35 = \qty{184.08}{Ma}$: los \qty{184.2}{Ma} de la tabla quedan
fuera, por \qty{0.12}{Myr}. (b) $183.73 - 182.77 = \qty{0.96}{Myr}$. (c) $6/0.96 =
\qty{6.25}{m}$ por millón de años. (d) Una biozona se define por fósiles, cuyas apariciones y
desapariciones son casi simultáneas en áreas extensas; el espesor de roca depositado durante
ese tiempo depende de la velocidad local de acumulación.
\end{solution}

\begin{solution}{exo:b2:palaeoenvironments:7}
(a) $0.49\,e^{-1046/3704} = 0.369$ y $0.49\,e^{-1166/3704} = 0.358$; media $0.364$. (b) $120
\times (1 - 0.364) = \qty{76.3}{m}$. (c) $76.3/(1 - 0.49) = \qty{149.6}{m}$. (d) El factor
supone que toda la capa restaurada tiene la porosidad $\varphi_0$ de la superficie; pero una capa
de \qty{150}{m} de espesor en superficie es menos porosa en su base que en su techo, de modo que los mismos granos
ocupan menos de \qty{149.6}{m}.
\end{solution}

\begin{solution}{exo:b2:palaeoenvironments:8}
(a) Un volumen unidad contiene $\varphi$ de agua y $1 - \varphi$ de granos: $\rho = \varphi
\rho_{\text{agua}} + (1 - \varphi)\rho_{\text{grano}}$, lo que da la fórmula; $(2710 -
2573)/(2710 - 1000) = 0.080$. (b) $0.51\,e^{-956/4545} = 0.413$. (c) La cementación: la calcita
precipitada en los poros, que sufren los carbonatos además de la compactación, elimina mucha más porosidad
que el enterramiento por sí solo.
\end{solution}

\begin{solution}{exo:b2:palaeoenvironments:9}
(a) Turbiditas de mar profundo (\omterm{def:b2:palaeoenvironments:ichnofacies}{icnofacies} de Nereites); un fondo jurásico de la plataforma inferior o del talud
(Zoophycos); un fondo de plataforma externa entre los niveles de base del oleaje (Cruziana); una playa sumergida (Skolithos); una
ciénaga en tierra firme. (b) El agua se hizo más somera hacia arriba: la costa se desplazó hacia el mar. (c) Un descenso del nivel del mar,
o un aporte de sedimento más rápido que la subsidencia, como delante de un delta que crece; un descenso
se encontraría en la misma época en otras cuencas, un delta sería local y su sedimento llegaría
de un solo lado. (d) En una sucesión concordante, los ambientes vecinos se suceden; un
salto sobre los intermedios indica una laguna temporal o una falla.
\end{solution}

\begin{solution}{exo:b2:palaeoenvironments:10}
(a) El tiempo hasta el enterramiento es $h/s$ = 100, 10 y \qty{1}{yr}. Cuerpo blando: $e^{-100} \approx
\num{4e-44}$, $e^{-10} = \num{4.5e-5}$, $e^{-1} = 0.37$. Concha ($k = \qty{0.01}{yr^{-1}}$):
$e^{-1} = 0.37$, $e^{-0.1} = 0.90$, $e^{-0.01} = 0.99$. (b) $s = kh/\ln 100 = 0.1/4.605 =
\qty{0.022}{m.yr^{-1}}$, \qty{22}{mm} al año. (c) Difícilmente, durante meses; pero una capa de tormenta, una
turbidita o una caída de ceniza entierran el fondo en horas. (d) La velocidad de sedimentación constante: el
tiempo hasta el enterramiento es entonces mucho más corto que $h/s$ con la velocidad media, y $P = e^{-kT}$, con $T$
el tiempo real hasta el enterramiento, tiende a uno cuando $T$ es mucho menor que $1/k$.
\end{solution}

\begin{solution}{exo:b2:palaeoenvironments:11}
(a) Una capa delgada $\mathrm dz$ a la profundidad $z$ contiene los granos $(1 - \varphi(z))\,\mathrm dz$;
en superficie ocupan $\mathrm dz_0$ con $(1 - \varphi_0)\,\mathrm dz_0 = (1 -
\varphi(z))\,\mathrm dz$. (b) $\varphi = 0.63\,e^{-816/1960} = 0.415$; $(1 - 0.415)/0.37 =
1.581$, así que \qty{158.1}{m}. (c) \qty{9.5}{m} de más: como en el \cref{exo:b2:palaeoenvironments:7},
la capa restaurada no es uniformemente tan porosa como su superficie. (d) La velocidad de granos,
$h(1-\varphi)/\Delta t$, solo usa el espesor y la porosidad actuales, ambos medidos.
\end{solution}

\begin{solution}{exo:b2:palaeoenvironments:12}
(a) $9.5/0.165 = 57.6$ horizontes. (b) $180/57.6 = \qty{3.1}{m}$ de caliza hoy, $338/57.6 =
\qty{5.9}{m}$ de sedimento fresco. (c) Como mucho \qty{1.5}{m}: $180/1.5 = 120$ intervalos, 121
muestras. (d) Interrupciones de la sedimentación, capas sin ammonites, y horizontes condensados en
unos pocos centímetros donde el sedimento llegaba despacio.
\end{solution}

\begin{solution}{pb:b2:palaeoenvironments:1}
\textbf{1.} $866 \sin \ang{12} = 866 \times 0.2079 = \qty{180.0}{m}$.
\textbf{2.} $481 \times 0.2079 = \qty{100.0}{m}$.
\textbf{3.} La arenisca, por superposición: los flancos buzan solo \ang{12}, así que las capas no están
invertidas y la más baja es la más antigua.
\textbf{4.} $866 - 766 = \qty{100}{m}$ y $1046 - 866 = \qty{180}{m}$.
\textbf{5.} Al menos \qty{766}{m} allí; su techo original está erosionado, así que su espesor completo se
desconoce.
\textbf{6.} 0.70 en superficie; $0.70\,e^{-766/1408} = 0.406$ a \qty{766}{m}.
\textbf{7.} Son grupos estenohalinos: el agua tenía la salinidad del mar abierto.
\textbf{8.} Una \omterm{def:b2:palaeoenvironments:ichnofacies}{icnofacies} de Cruziana: la plataforma externa, entre los niveles de base del oleaje de buen tiempo y de
tormenta.
\textbf{9.} Tormentas: sus olas llegaban al fondo, rompían las conchas y las depositaban en capas
de base erosiva.
\textbf{10.} El agua del fondo era pobre en oxígeno: pocos excavadores, láminas conservadas, formación de pirita.
\textbf{11.} La \omterm{def:b2:palaeoenvironments:ichnofacies}{icnofacies} de Zoophycos, en el Mesozoico la plataforma inferior o el talud, por debajo del
nivel de base del oleaje de tormenta; con oxígeno suficiente a veces para los excavadores.
\textbf{12.} No: los ammonites nadaban en el agua de encima; sus conchas se hundieron o derivaron hasta el fondo
tras la muerte y forman una \omterm{def:b2:palaeoenvironments:assemblages}{asociación de muerte}.
\textbf{13.} $184.2 - 174.7 = \qty{9.5}{Myr}$.
\textbf{14.} $180/9.5 = \qty{18.9}{m}$ por millón de años.
\textbf{15.} $(2710 - 2573)/(2710 - 1000) = 0.080$.
\textbf{16.} $180 \times (1 - 0.080)/(1 - 0.51) = \qty{338}{m}$.
\textbf{17.} $338/9.5 = \qty{35.6}{m}$ por millón de años.
\textbf{18.} $180 \times 0.92/9.5 = \qty{17.4}{m}$ por millón de años.
\textbf{19.} La lutita daba 11.3 y \qty{4.4}{m} por millón de años: la velocidad de granos se redujo
unas cuatro veces. El carbonato del bajío y de la plataforma dejó de llegar a medida que el fondo
se hacía más profundo, y solo lo alcanzaba una lenta lluvia de fango.
\textbf{20.} Avanzó: ríos, después un bajío, un fondo de plataforma externa, un fondo por debajo del nivel de base del oleaje de tormenta,
y después el mar abierto de la creta. La columna registra una transgresión.
\textbf{21.} $57 - 41 = \qty{16}{m}$ durante la caliza, $86 - 57 = \qty{29}{m}$ durante la lutita,
\qty{45}{m} en total.
\textbf{22.} $(60 - 5) + 338 + 148.6 = \qty{542}{m}$, en $184.2 - 161.5 = \qty{22.7}{Myr}$:
\qty{23.9}{m} por millón de años.
\textbf{23.} $45/542 = 0.083$: alrededor del \qty{8}{\percent}, y menos si la carta de ciclos exagera
la amplitud. El resto procedió de la subsidencia del fondo.
\textbf{24.} El fango calcáreo de la caliza de Varden se acumuló a unos
\textbf{\qty[text-series-to-math=true]{36}{m}} por millón de años cuando se depositó.
\end{solution}
